

\begin{enumerate}[label=\thesubsection.\arabic*, ref=\thesubsection.\theenumi]

\item
\begin{align}
\mydet{a & b & c\\d & e & f\\g & h & i} = a\brak{ei - fh} - b\brak{di - fg} + c\brak{dh - eg}
\label{eq:EE26BTECH11211-Prob_1}
\end{align}

\item 
King's Rule

\begin{align}
\int_a^b f\brak{x}\,dx
&= \int_a^b f\brak{a+b-x}\,dx
\label{eq:EE26BTECH11211-Prob_2.1}
\end{align}

For the interval $[0,a]$,

\begin{align}
\int_0^a f\brak{x}\,dx
&= \int_0^a f\brak{a-x}\,dx
\label{eq:EE26BTECH11211-Prob_2.2}
\end{align}

Hence,

\begin{align}
2\int_0^a f\brak{x}\,dx
&= \int_0^a \sbrak{f\brak{x}+f\brak{a-x}}\,dx
\label{eq:EE26BTECH11211-Prob_2.3}
\end{align}

\item 
Integration by Parts

\begin{align}
\int u\,dv &= uv-\int v\,du
\label{eq:EE26BTECH11211-Prob_3.1}
\end{align}

For definite integrals,

\begin{align}
\int_a^b u\,dv
&= \sbrak{uv}_a^b-\int_a^b v\,du
\label{eq:EE26BTECH11211-Prob_3.2}
\end{align}

\item 
Integration by Substitution

\begin{align}
u &= g\brak{x}, \qquad du=g'\brak{x}\,dx
\end{align}

\begin{align}
\int f\brak{g\brak{x}}g'\brak{x}\,dx
&= \int f\brak{u}\,du
\label{eq:EE26BTECH11211-Prob_4.1}
\end{align}

For definite integrals,

\begin{align}
\int_a^b f\brak{g\brak{x}}g'\brak{x}\,dx
&= \int_{g\brak{a}}^{g\brak{b}} f\brak{u}\,du
\label{eq:EE26BTECH11211-Prob_4.2}
\end{align}

\begin{align}
\int \frac{f'\brak{x}}{f\brak{x}}\,dx
&= \ln\abs{f\brak{x}}+C
\label{eq:EE26BTECH11211-Prob_4.3}
\end{align}

\item 
Partial Fractions

For distinct constants $a$ and $b$,

\begin{align}
\frac{1}{\brak{x-a}\brak{x-b}}
&= \frac{1}{a-b}
\left[
\frac{1}{x-a}-\frac{1}{x-b}
\right]
\end{align}

\begin{align}
\int\frac{dx}{\brak{x-a}\brak{x-b}}
&= \frac{1}{a-b}
\ln\abs{\frac{x-a}{x-b}}+C
\end{align}

\begin{align}
\frac{1}{\brak{x-a}\brak{x-b}}
&= \frac{A}{x-a}+\frac{B}{x-b}
\label{eq:EE26BTECH11211-Prob_5.1}
\end{align}

where

\begin{align}
A &= \frac{1}{a-b},
& B &= -\frac{1}{a-b}
\end{align}

\item
Finding Area

Area under the curve $y=f\brak{x}$, where $f\brak{x}\geq 0$,

\begin{align}
A &= \int_a^b f\brak{x}\,dx\label{eq:EE26BTECH11211-Prob_6.1}
\end{align}

Area between a curve and the $x$-axis,

\begin{align}
A &= \int_a^b \abs{f\brak{x}}\,dx
\end{align}

If the curve crosses the $x$-axis, split the integral at its points of intersection.

Area between two curves, where $f\brak{x}\geq g\brak{x}$,

\begin{align}
A &= \int_a^b \sbrak{f\brak{x}-g\brak{x}}\,dx
\end{align}

In general,

\begin{align}
A &= \int_a^b \abs{f\brak{x}-g\brak{x}}\,dx
\end{align}

Area with respect to $y$, for the curve $x=f\brak{y}$ and the $y$-axis,

\begin{align}
A &= \int_c^d \abs{f\brak{y}}\,dy
\end{align}

Area between two curves with respect to $y$,

\begin{align}
A &= \int_c^d
\sbrak{x_{\text{right}}\brak{y}-x_{\text{left}}\brak{y}}\,dy
\end{align}

\item 
Let $\vec{a}$ and $\vec{b}$ are two vectors then 
\begin{align}
\vec{a}.\vec{b} = |a||b|cos\theta
\label{eq:EE26BTECH11211-Prob_7.1}
\end{align}
where $\theta$ is the angle between two vectors.

\item
\begin{align}
\dfrac{d}{dx}x^n = nx^{n-1}
\end{align}

\item
The conditional probability of event $S$ given event $E$
is defined as
\begin{align}
P\brak{S\mid E}=\frac{P\brak{S\cap E}}{P\brak{E}}.
\label{eq:EE26BTECH11211-Prob_9.1}
\end{align}

\item
\begin{align}
\frac{d}{dx}\left(\tan^{-1}u\right)
= \frac{1}{1+u^2}\frac{du}{dx}.
\label{eq:EE26BTECH11211_Prob_10.1}
\end{align}

\item
Let $\vec{A} = x_1\vec{i} + y_1\vec{j} + z_1\vec{k}$ and $\vec{B}= x_2\vec{i} + y_2\vec{j} + z_2\vec{k}$ be two vectors then the vector perpendicular to both the vectors is 
\begin{align}
\vec{A}\times\vec{B}
=
\myvec{
a_2b_3 - a_3b_2\\
a_3b_1 - a_1b_3\\
a_1b_2 - a_2b_1
}\label{eq:EE26BTECH11211_Prob_11.1}
\end{align}

\item
Let $\vec{A}$ be a vector then the unit vector in the direction of vector $\vec{A}$
is $\dfrac{\vec{A}}{\abs{A}}$\label{eq:EE26BTECH11211_Prob_12.1}

\item
$\dfrac{de^{ax}}{dx} = e^{ax}a$ \label{eq:EE26BTECH11211_Prob_13}

\item 
Linear first-order differential equation \label{eq:EE26BTECH11211_Prob_14}

For
\begin{align}
    \frac{dy}{dx}+P(x)y=Q(x)
\end{align}

The integrating factor is
\begin{align}
    \mathrm{I.F.}=e^{\int P(x)\,dx}
\end{align}

The solution is
\begin{align}
    y\,e^{\int P(x)\,dx}
    =\int Q(x)e^{\int P(x)\,dx}\,dx+C
\end{align}

\item
$\dfrac{d}{dx}sin^{-1}x = \dfrac{1}{\sqrt1-x^2}$\label{eq:EE26BTECH11211_Prob_15}

\item
 Bayes' theorem ,
\begin{align}
P\brak{E\mid A}
=\frac{P\brak{E}P\brak{A\mid E}}{P\brak{A}}\label{eq:EE26BTECH11211-Prob_16}.
\end{align}

\item
The following Laplace transform formulae are used in the solutions.

\begin{align}
\mathcal{L}\{f(t)\}
&=F(s)=\int_0^\infty e^{-st}f(t)\,dt,
\label{eq:EE26BTECH11211_Prob_1.1.17.1}
\\
\mathcal{L}\left\{\frac{df}{dt}\right\}
&=sF(s)-f(0),
\label{eq:EE26BTECH11211_Prob_1.1.17.2}
\\
\mathcal{L}\{t f(t)\}
&=-\frac{dF(s)}{ds},
\label{eq:EE26BTECH11211_Prob_1.1.17.3}
\\
\mathcal{L}\{1\}
&=\frac{1}{s},
\label{eq:laplace-one}
\\
\mathcal{L}\{t^n\}
&=\frac{n!}{s^{n+1}},\qquad n>-1,
\label{eq:EE26BTECH11211_Prob_1.1.17.5}
\\
E_1(s)&=\int_s^\infty \frac{e^{-u}}{u}\,du,
\label{eq:exponential-integral}
\\
\mathcal{L}\left\{\frac{1}{1+t}\right\}
&=e^sE_1(s),
\label{eq:EE26BTECH11211_Prob_1.1.17.7}
\\
\mathcal{L}\{\ln(1+t)\}
&=\frac{e^sE_1(s)}{s},
\label{eq:EE26BTECH11211_Prob_1.1.17.8}
\\
\mathcal{L}^{-1}\left\{\frac{1}{s^{n+1}}\right\}
&=\frac{t^n}{n!}.
\label{eq:EE26BTECH11211_Prob_1.1.17.9}
\end{align}
\end{enumerate}

